\documentclass[preprint,a4paper]{article}
\usepackage[margin=1in]{geometry}
\usepackage{amsmath,amssymb,amsthm}
\usepackage{graphicx}
\usepackage{booktabs}
\usepackage{hyperref}
\usepackage{xcolor}
\usepackage{listings}
\usepackage{algorithm}
\usepackage{algorithmic}

% Theorem environments
\newtheorem{theorem}{Theorem}
\newtheorem{lemma}{Lemma}
\newtheorem{definition}{Definition}
\newtheorem{claim}{Claim}
\newtheorem{remark}{Remark}

% Title
\title{
  \textbf{QuantRot-PQC: HD$^3$ Rotation Reduces\\
  KV-Cache Coefficient Dynamic Range for\\
  Post-Quantum Secure LLM Inference}
}

\author{
  Venkata Vijaya Kumari D.\\
  Independent Researcher\\
  \texttt{vijayaqa.rv7@gmail.com}\\
  \url{https://github.com/vijayarjun7/QuantRot-PQC}
}

\date{July 2026}

\begin{document}

\maketitle

\begin{abstract}
LLM inference pipelines transmitting KV-cache between
nodes face harvest-now-decrypt-later (HNDL) attacks:
adversaries record encrypted traffic today for
retroactive decryption by future quantum computers.
CRYSTALS-Kyber provides post-quantum protection but
adds NTT polynomial arithmetic overhead at inference
time. We propose \textbf{QuantRot-PQC}, a
pre-encryption transformation that applies HD$^3$
random unitary rotation to KV-cache vectors prior to
Kyber polynomial encoding. By redistributing vector
energy across dimensions, the rotation reduces the
centered L$^\infty$ norm of input polynomials,
potentially reducing coefficient dynamic range before
lattice encoding.

We demonstrate that HD$^3$ rotation reduces
L$^\infty$ norm by \textbf{40.1\%} on 500 real
KV vectors from Llama-3.2-1B --- double the
$\geq 20\%$ threshold --- while kurtosis drops
from $6.15$ to $-0.09$, confirming near-Gaussian
concentration. A null result on synthetic isotropic
Gaussian vectors ($0.2\%$ reduction) serves as an
experimental control: rotation cannot concentrate
what is already uniformly distributed.
Pipeline latency in Python is dominated by
preprocessing overhead at $d = 64$; a compiled
HD$^3$ kernel is required for production deployment.
The pre-rotation is public, invertible,
input-independent, and operates in the plaintext
domain prior to Kyber's algebraic operations ---
it does not alter the Module-LWE problem instance
and is expected to inherit Kyber's IND-CCA2 security
under standard assumptions. In systems already using
TurboQuant KV-cache compression, QuantRot-PQC adds
post-quantum security with zero additional rotation
cost. Code and experiments are publicly available.
\end{abstract}

\textbf{Keywords:} post-quantum cryptography,
CRYSTALS-Kyber, KV-cache, LLM inference,
HD$^3$ rotation, t-designs, Module-LWE



%% ─────────────────────────────────────────────
\section{Introduction}
\label{sec:intro}
Large language model (LLM) inference pipelines that
transmit key-value (KV) cache between nodes ---
across multi-GPU systems, edge-cloud boundaries, or
federated inference networks --- face an emerging
security threat: \emph{harvest-now-decrypt-later}
(HNDL) attacks. Adversaries record encrypted
inference traffic today and decrypt it retroactively
when fault-tolerant quantum computers arrive,
exposing model reasoning, retrieved context, and
user data. CRYSTALS-Kyber~\cite{kyber2021}, the
NIST-standardized post-quantum KEM, provides
quantum-resistant protection for such transmissions.
However, Kyber's NTT polynomial arithmetic and
$\approx 2\times$ ciphertext expansion add latency
and memory overhead that must be minimized for
production deployment.

\textbf{The gap.}
Kyber is designed for arbitrary data and does not
exploit the mathematical structure of its input.
Real LLM KV-cache vectors, however, have a specific
and exploitable property: after layer normalization,
they exhibit \emph{heavy-tailed} coordinate
distributions with significant outlier dimensions.
We confirm this empirically on Llama-3.2-1B,
finding mean kurtosis $6.15$ --- far above the
kurtosis $\approx 0$ of isotropic Gaussian vectors.
When such vectors are encoded directly into Kyber's
polynomial ring, large outlier coefficients may
increase implementation-specific overhead in
NTT arithmetic. No existing work exploits this
structure to improve Kyber encoding efficiency
for LLM inference.

\textbf{Our contribution.}
We propose \textbf{QuantRot-PQC}, a pre-encryption
transformation that applies an HD$^3$ random unitary
rotation to KV-cache vectors before CRYSTALS-Kyber
polynomial encoding. By redistributing vector energy
across dimensions, the rotation converts heavy-tailed
distributions to near-Gaussian, reducing the
centered L$^\infty$ norm and enabling more uniform
polynomial coefficient distributions.

Specifically, we make the following contributions:

\begin{enumerate}

    \item \textbf{QuantRot-PQC construction}
    (Section~\ref{sec:method}): A formal pipeline
    mapping KV vectors $x \in \mathbb{R}^d$ through
    HD$^3$ rotation, scaling to centered $\mathbb{Z}_q$
    coefficients, and Kyber-512 encapsulation.
    The transformation is public, invertible,
    input-independent, and operates entirely in the
    plaintext domain prior to Kyber's algebraic
    operations.

    \item \textbf{Security argument}
    (Section~\ref{sec:security}): We argue that
    QuantRot-PQC is expected to inherit Kyber's
    IND-CCA2 security under standard Module-LWE
    assumptions, as the pre-rotation does not alter
    the Module-LWE problem instance
    $(A, A\mathbf{s}+\mathbf{e})$.

    \item \textbf{Empirical validation}
    (Section~\ref{sec:experiments}): We extract
    500 real KV vectors from Llama-3.2-1B via
    \texttt{past\_key\_values} and demonstrate:
    \begin{itemize}
        \item HD$^3$ rotation reduces L$^\infty$
        norm by $\mathbf{40.1\%}$ on real KV
        vectors --- double the $\geq 20\%$
        threshold --- while kurtosis drops from
        $6.15$ to $-0.09$.
        \item Synthetic isotropic Gaussian vectors
        yield a null result ($0.2\%$ reduction),
        as expected from rotation invariance.
        This serves as an experimental control.
        \item Pipeline latency in Python is
        dominated by preprocessing overhead at
        $d = 64$; a compiled HD$^3$ kernel is
        required for production deployment.
    \end{itemize}

\end{enumerate}

\textbf{Relationship to TurboQuant.}
TurboQuant~\cite{turboquant2025} demonstrated
that HD$^3$ rotation enables near-Shannon-optimal
KV-cache compression by producing concentrated
coordinate distributions. QuantRot-PQC applies
the same rotation for a distinct purpose:
reducing polynomial coefficient dynamic range
before Kyber encryption rather than before
scalar quantization. In systems already using
TurboQuant, QuantRot-PQC adds post-quantum
security with zero additional rotation cost ---
the rotation is shared between compression and
encryption pipelines.

\textbf{Paper organization.}
Section~\ref{sec:background} covers background
on CRYSTALS-Kyber, KV-cache structure, and
random rotation theory.
Section~\ref{sec:method} presents the
QuantRot-PQC construction and security argument.
Section~\ref{sec:experiments} reports experimental
results.
Section~\ref{sec:discussion} discusses limitations
and future work.
Section~\ref{sec:conclusion} concludes.

%% ─────────────────────────────────────────────
\section{Background}
\label{sec:background}
\subsection{CRYSTALS-Kyber}
\label{sec:kyber}

CRYSTALS-Kyber~\cite{kyber2021} is a lattice-based key
encapsulation mechanism (KEM) standardized by NIST for
post-quantum cryptography. Kyber operates over the
polynomial ring $\mathcal{R}_q = \mathbb{Z}_q[X]/(X^n+1)$
where $n = 256$ and $q = 3329$. The scheme is based on
the hardness of the \emph{Module Learning With Errors}
(Module-LWE) problem over a $k \times k$ matrix of
polynomials, where $k = 2$ for Kyber-512 (NIST Level 1).

\textbf{Key generation} (Algorithm 4,~\cite{kyber2021})
samples a public matrix $A \in \mathcal{R}_q^{k \times k}$,
secret vector $\mathbf{s} \in \mathcal{R}_q^k$, and
error vector $\mathbf{e} \in \mathcal{R}_q^k$, producing
public key $\mathbf{t} = A\mathbf{s} + \mathbf{e}$.

\textbf{Encryption} (Algorithm 5,~\cite{kyber2021})
takes a plaintext message $m \in \{0,1\}^{256}$ and
produces ciphertext $(\mathbf{u}, v)$ via:
\[
    \mathbf{u} = A^\top \mathbf{r} + \mathbf{e}_1,
    \quad
    v = \mathbf{t}^\top \mathbf{r} + e_2 +
    \mathsf{Decompress}_q(m, 1)
\]
where $\mathbf{r}, \mathbf{e}_1, e_2$ are freshly sampled
noise terms independent of $m$. The plaintext $m$ enters
the pipeline at this step --- specifically as the argument
to $\mathsf{Decompress}_q$ in Algorithm 5, Line 20.
This is the exact insertion point of QuantRot-PQC.

\textbf{Security:} Kyber-512 achieves NIST Level 1
security (equivalent to AES-128), with 118 bits of
classical hardness and 107 bits of quantum hardness,
estimated using the framework of~\cite{albrecht2015}.
The IND-CCA2 security proof (Theorem~2,~\cite{kyber2021})
holds for arbitrary message distributions --- a key
property exploited by QuantRot-PQC.
Decryption failure probability is $\delta = 2^{-139}$,
governed solely by noise terms independent of $m$.

\textbf{Performance:} Kyber-512 encapsulation requires
approximately 45,200 AVX2 cycles ($\approx 6~\mu$s at
3.5~GHz) on commodity x86 hardware. We confirm
$6.25~\mu$s mean encapsulation latency in our baseline
experiments (Section~\ref{sec:experiments}).

\subsection{KV-Cache Structure in LLM Inference}
\label{sec:kvcache}

Transformer attention mechanisms compute key and value
projections for each input token, caching them for
subsequent token generation. For a model with $H$
attention heads and head dimension $d$, each token
produces $H$ key vectors and $H$ value vectors, each
in $\mathbb{R}^d$.

After layer normalization, KV vectors are approximately
unit-norm. However, unlike synthetic Gaussian vectors,
real KV vectors exhibit \emph{heavy-tailed} coordinate
distributions with significant outlier dimensions --- a
known property of transformer activations
\cite{turboquant2025}.

We confirm this empirically on Llama-3.2-1B: KV vectors
extracted via \texttt{past\_key\_values} have mean
L$^\infty$ norm $0.5407$ and kurtosis $6.15$, compared
to kurtosis $\approx 0$ for isotropic Gaussian vectors.
This non-Gaussian structure is the property exploited
by QuantRot-PQC.

\textbf{Security threat:} LLM inference pipelines that
transmit KV-cache between nodes (multi-GPU, edge-cloud,
federated inference) are vulnerable to
\emph{harvest-now-decrypt-later} (HNDL) attacks.
Adversaries record encrypted inference traffic today
and decrypt it retroactively when fault-tolerant quantum
computers arrive. CRYSTALS-Kyber provides quantum-resistant
protection but adds overhead that must be minimized for
production deployment.

\subsection{Random Rotation and t-Designs}
\label{sec:tdesigns}

A \emph{unitary t-design} is an ensemble of unitaries
whose first $t$ moments match those of the Haar
(uniformly random) measure over unitary matrices.
For our purposes, an approximate 2-design is sufficient:
it ensures that coordinates of rotated vectors are
approximately independent with variance $\approx 1/d$.

Brandão et al.~\cite{brandao2016} showed that local
random quantum circuits of depth $O(t^{10} n^2)$ form
approximate unitary t-designs. Harrow and
Mehraban~\cite{harrow2018} improved this to depth
$O(\text{poly}(t) \cdot n^{1/D})$ for D-dimensional
lattice architectures.

The HD$^3$ construction --- three applications of
Hadamard transform composed with random $\pm 1$ diagonal
matrices --- is a classical structured approximation
consistent with these results, achieving the
concentration property required for our analysis at
$O(d \log d)$ cost per vector.

\textbf{Connection to TurboQuant:}
TurboQuant~\cite{turboquant2025} demonstrated that random
rotation converts arbitrary KV vectors to approximately
isotropic distributions, enabling near-Shannon-optimal
scalar quantization. Specifically, after rotation,
coordinates follow a Beta$(1/2, (d-1)/2)$ distribution
that converges to $\mathcal{N}(0, 1/d)$ for large $d$.
QuantRot-PQC applies the same rotation for a distinct
purpose: reducing $\ell^\infty$ norm before Kyber
polynomial encoding, rather than for compression.

%% ─────────────────────────────────────────────
\section{QuantRot-PQC}
\label{sec:method}

We propose QuantRot-PQC, a pre-encryption transformation
that applies an HD$^3$ random unitary rotation to LLM
KV-cache vectors prior to CRYSTALS-Kyber polynomial
encoding. The transformation exploits the heavy-tailed
non-Gaussian structure of real KV vectors to reduce
centered coefficient dynamic range before lattice encoding.

\subsection{KV Vector to Polynomial Mapping}
\label{sec:mapping}

Let $x \in \mathbb{R}^d$ be a KV-cache vector extracted
from a transformer attention layer, unit-normalized after
layer normalization ($\|x\|_2 = 1$). The QuantRot-PQC
encoding pipeline proceeds as follows:

\begin{enumerate}
    \item \textbf{Rotation:} Apply HD$^3$ rotation $U$:
    \[
        x' = U \cdot x, \quad \|x'\|_2 = 1
    \]

    \item \textbf{Scaling to $\mathbb{Z}_q$:} Map to
    centered polynomial coefficients:
    \[
        c_i = \text{round}(\alpha \cdot x'_i) \bmod q,
        \quad c_i \in \left[-\frac{q}{2}, \frac{q}{2}\right]
    \]
    where $q = 3329$ (Kyber modulus) and scale factor
    $\alpha \approx 5235$.

    \item \textbf{Kyber encoding:} Encode $\{c_i\}$ as
    plaintext message $m \in \mathcal{R}_q$ and encapsulate:
    \[
        \mathsf{ct} = \mathsf{Kyber.Enc}(\mathsf{pk}, m)
    \]
\end{enumerate}

The scale factor $\alpha$ is derived from the theoretical
$\ell^\infty$ bound after rotation. For unit-norm vectors
after HD$^3$ rotation, coordinates are sub-Gaussian with
variance proxy $\approx 1/d$. By the union bound over
$d$ coordinates:
\[
    \|x'\|_\infty \leq
    \sqrt{\frac{2\log(2d/\delta)}{d}}
    \quad \text{with probability } 1 - \delta
\]
For $d = 64$, $\delta = 0.001$, this gives
$\|x'\|_\infty \leq 0.6062$.

\textbf{Note on centered representation:}
Negative coordinates $x'_i < 0$ map to negative centered
coefficients, represented in standard $\mathbb{Z}_q$ as
$c_i + q$. The $\ell^\infty$ claim holds in the centered
representation $[-q/2, q/2]$, not the unsigned
representation $[0, q-1]$.

\subsection{HD$^3$ Rotation Construction}
\label{sec:hd3}

The HD$^3$ rotation $U$ is defined as:
\[
    U = H \cdot D_3 \cdot H \cdot D_2 \cdot H \cdot D_1
\]
where $H$ is the Walsh-Hadamard transform and each
$D_i = \text{diag}(r_1, \ldots, r_d)$ with
$r_j \sim \text{Uniform}\{-1, +1\}$ independently sampled.

\textbf{Properties:}
\begin{itemize}
    \item \textbf{Orthogonal:} $UU^\top = I$,
          preserving $\ell_2$ norm.
    \item \textbf{Efficient:} $O(d \log d)$ per vector.
    \item \textbf{Public:} Sign vectors $D_1, D_2, D_3$
          shared openly. Stored as $3d$ bits.
    \item \textbf{Invertible:}
          $(HD^3)^{-1} = (HD^3)^\top$.
\end{itemize}

\subsection{Security Argument}
\label{sec:security}

\begin{claim}
QuantRot-PQC is expected to inherit the IND-CCA2 security
of the underlying CRYSTALS-Kyber scheme under standard
Module-LWE assumptions.
\end{claim}

\textbf{Argument:}
The pre-rotation $U$ is public, invertible,
input-independent, and operates entirely in the plaintext
domain prior to Kyber's algebraic operations.
The Module-LWE problem instance
$(A, \mathbf{t} = A\mathbf{s} + \mathbf{e})$ is generated
independently of the plaintext message $m$.
Pre-rotation modifies only $m$ --- it does not alter
the distributions of $A$, $\mathbf{s}$, or $\mathbf{e}$.
Consequently, the construction is expected to inherit
Kyber's IND-CCA2 security guarantee
(Theorem~2,~\cite{kyber2021}).

Kyber's decryption failure probability
$\delta = 2^{-139}$ depends solely on noise polynomials
$\mathbf{e}_1, \mathbf{e}_2 \sim B_\eta$, sampled
independently of $m$. Pre-rotation does not affect
$\delta$.

\textbf{Formal reduction:} A complete IND-CCA2 security
reduction is left for future work.

\subsection{Effect on Coefficient Dynamic Range}
\label{sec:dynamic_range}

Real KV vectors have heavy-tailed coordinate
distributions --- a single outlier dimension dominates
the $\ell^\infty$ norm. HD$^3$ rotation redistributes
this energy across all $d$ dimensions, producing
concentrated near-isotropic coordinates.
The practical consequence is measured in
Section~\ref{sec:experiments}: $\ell^\infty$ norm reduces
by 40.1\% on real Llama-3.2-1B KV vectors.

\textbf{Important caveat:} The actual NTT reduction
frequency inside \texttt{liboqs} is not directly
measurable via the public API ---
\texttt{encap\_secret()} generates its own randomness
internally. The relationship between input coefficient
range and NTT implementation efficiency requires an
instrumented Kyber implementation and is left as
future work.

%% ─────────────────────────────────────────────
\section{Experiments}
\label{sec:experiments}
\subsection{Experimental Setup}
\label{sec:setup}

All experiments run on a MacBook Air (Apple M4,
24GB RAM, macOS 26.5.1) with Python 3.14.2.
Kyber-512 encapsulation uses \texttt{liboqs-python}
(Open Quantum Safe project~\cite{liboqs}), which
auto-builds the \texttt{liboqs} C library on first
import. Real KV vectors are extracted from
\texttt{meta-llama/Llama-3.2-1B}~\cite{llama2024}
via \texttt{past\_key\_values} using the HuggingFace
Transformers library. All experiments use
\texttt{numpy.random.seed(42)} for reproducibility.
Code is publicly available at:
\url{https://github.com/vijayarjun7/QuantRot-PQC}

\subsection{Baseline: Kyber-512 Latency and KV
Distribution}
\label{sec:baseline}

Table~\ref{tab:baseline} reports Kyber-512
encapsulation latency (100 trials) and KV vector
distribution statistics (1000 synthetic unit-norm
Gaussian vectors, $d = 256$) before any rotation.

\begin{table}[h]
\centering
\caption{Baseline measurements before rotation.}
\label{tab:baseline}
\begin{tabular}{lrr}
\toprule
\textbf{Metric} & \textbf{Synthetic ($d$=256)}
& \textbf{Real Llama-3.2-1B ($d$=64)} \\
\midrule
Kyber-512 latency (mean) & \multicolumn{2}{c}{6.25 $\mu$s} \\
Kyber-512 latency (std)  & \multicolumn{2}{c}{0.91 $\mu$s} \\
Ciphertext size          & \multicolumn{2}{c}{768 bytes (spec)} \\
\midrule
L$^\infty$ norm (mean) & 0.1907 & 0.5407 \\
Kurtosis (mean)        & $\approx 0.0$ & 6.15 \\
\% within $\pm q/2$    & 100\% & 100\% \\
\bottomrule
\end{tabular}
\end{table}

The contrast between synthetic and real KV vectors
is striking. Synthetic isotropic Gaussian vectors
have kurtosis $\approx 0$ and L$^\infty$ norm
$0.1907$, already close to the theoretical bound
of $0.3205$ for $d = 256$, $\delta = 0.001$.
Real Llama-3.2-1B KV vectors have kurtosis $6.15$
--- confirming heavy-tailed non-Gaussian structure
--- and L$^\infty$ norm $0.5407$, substantially
higher than the synthetic baseline.

\subsection{H1: L$^\infty$ Norm Reduction}
\label{sec:h1}

\textbf{Hypothesis H1:} HD$^3$ rotation reduces the
empirical L$^\infty$ norm of KV vectors by $\geq 20\%$
compared to unrotated vectors.

We test H1 on two datasets: synthetic isotropic
Gaussian vectors and real KV vectors from
Llama-3.2-1B. Table~\ref{tab:h1} reports results
for both.

\begin{table}[h]
\centering
\caption{H1 results: L$^\infty$ norm before and
after HD$^3$ rotation.}
\label{tab:h1}
\begin{tabular}{lrrrr}
\toprule
\textbf{Dataset} & \textbf{Before} & \textbf{After}
& \textbf{Reduction} & \textbf{H1 Status} \\
\midrule
Synthetic ($d$=256, $n$=1000) &
0.1907 & 0.1904 & 0.2\% &
\textbf{Not Supported} \\
Real Llama-3.2-1B ($d$=64, $n$=500) &
0.5407 & 0.3239 & \textbf{40.1\%} &
\textbf{Supported} \\
\bottomrule
\end{tabular}
\end{table}

\textbf{Synthetic null result (expected):}
Isotropic Gaussian vectors are rotation-invariant
in distribution --- an orthogonal transform cannot
reduce the L$^\infty$ norm of vectors already
uniformly distributed across dimensions. The $0.2\%$
null result confirms theoretical expectations and
serves as an experimental control.

\textbf{Real KV result (main contribution):}
Real Llama-3.2-1B KV vectors yield a $40.1\%$
L$^\infty$ reduction --- double the $\geq 20\%$
threshold. Kurtosis drops from $6.15$ to $-0.09$,
confirming near-Gaussian concentration after rotation.
Figure~\ref{fig:linf} illustrates the distribution
shift.

\begin{figure}[h]
\centering
\includegraphics[width=0.85\textwidth]
{linf_real_comparison.png}
\caption{L$^\infty$ norm distribution and kurtosis
before vs.\ after HD$^3$ rotation on real
Llama-3.2-1B KV vectors ($n=500$, $d=64$).
Left: L$^\infty$ norm shifts from mean $0.5407$
to $0.3239$ ($40.1\%$ reduction).
Right: kurtosis drops from $6.15$ to $-0.09$,
confirming near-Gaussian concentration.}
\label{fig:linf}
\end{figure}

\textbf{Invertibility:}
The HD$^3$ rotation is numerically invertible with
maximum reconstruction error $2.98 \times 10^{-8}$
across all trials, confirming that the receiver can
exactly recover the original KV vector after
decapsulation.

\subsection{H2: Pipeline Latency}
\label{sec:h2}

\textbf{Hypothesis H2 (revised):} The full
QuantRot-PQC pipeline (HD$^3$ rotate + scale to
$\mathbb{Z}_q$ + encode + \texttt{encap\_secret})
adds acceptable overhead compared to vanilla
\texttt{encap\_secret} alone.

\textbf{Important caveat:} \texttt{liboqs}'s
\texttt{encap\_secret()} is a black box that
generates its own randomness internally and does
not accept structured polynomial input. We therefore
measure total pipeline latency, not internal NTT
reduction pressure. The latter requires an
instrumented Kyber implementation and is left as
future work.

Note that this benchmark tests $d = 64$ (the actual
Llama-3.2-1B head dimension) rather than $d = 256$
as in the original H2 statement; wall-clock $\mu$s
are reported via \texttt{perf\_counter\_ns()} rather
than hardware cycle counts, so the cited AVX2 cycle
figure and these latencies are not directly comparable.

Table~\ref{tab:h2} reports latency for three
HD$^3$ implementations at $d = 64$.

\begin{table}[h]
\centering
\caption{H2 pipeline latency comparison at $d=64$
(500 trials each).}
\label{tab:h2}
\begin{tabular}{lrrr}
\toprule
\textbf{Implementation} & \textbf{Preprocessing}
& \textbf{Total} & \textbf{vs Vanilla} \\
\midrule
Vanilla Kyber-512 & --- & 6.25 $\mu$s & baseline \\
Pure Python FWHT  & 110.89 $\mu$s & 117.94 $\mu$s
& $-1954\%$ \\
Matmul (BLAS)     & 32.01 $\mu$s  & 40.46 $\mu$s
& $-560\%$ \\
Reshape FWHT & 46.41 $\mu$s & 53.22 $\mu$s & $-752\%$ \\

\bottomrule
\end{tabular}
\end{table}

\textbf{H2 verdict: Not Supported} in the current
Python implementation. All three HD$^3$ variants
add substantial preprocessing overhead at $d = 64$.

\textbf{Root cause:} At $d = 64$, Kyber-512
compiled C ($\approx 6~\mu$s) is faster than any
Python-side preprocessing. The matmul approach
(32~$\mu$s) outperforms the asymptotically superior
reshape FWHT (48~$\mu$s) due to a single BLAS call
vs.\ 18 numpy reshape/slice operations --- numpy
per-call overhead dominates at this small dimension.

This is an \emph{implementation confound}, not
evidence against the algorithm. A compiled HD$^3$
kernel (C extension or CUDA) would reduce
preprocessing to $< 1~\mu$s, bringing the total
pipeline within $\approx 2\times$ of vanilla
Kyber-512. Figure~\ref{fig:h2} illustrates the
latency distributions.

\begin{figure}[h]
\centering
\includegraphics[width=0.75\textwidth]
{latency_comparison.png}
\caption{Latency distributions: vanilla Kyber-512
(mean $6.25~\mu$s, compiled C) vs.\ QuantRot-PQC
pipeline (mean $53.22~\mu$s, reshape FWHT + Kyber).
The gap reflects Python numpy overhead at $d = 64$,
not an inherent algorithmic cost.}
\label{fig:h2}
\end{figure}

%% ─────────────────────────────────────────────
\section{Discussion}
\label{sec:discussion}
\subsection{Why Real KV Vectors Differ From Synthetic}
\label{sec:synthetic_vs_real}

The contrast between the synthetic null result ($0.2\%$
reduction) and the real KV result ($40.1\%$ reduction)
is not a contradiction --- it is a confirmation of the
theoretical mechanism.

Synthetic isotropic Gaussian vectors have kurtosis
$\approx 0$ and are rotation-invariant in distribution.
An orthogonal transform cannot reduce the L$^\infty$
norm of vectors whose coordinates are already
identically and independently distributed.
The null result on synthetic data is therefore
expected and serves as an experimental control.

Real transformer KV vectors exhibit heavy-tailed
coordinate distributions (kurtosis $6.15$ on
Llama-3.2-1B). A small number of dimensions carry
disproportionate energy, inflating the L$^\infty$
norm. HD$^3$ rotation redistributes this energy
across all $d$ dimensions, converting heavy-tailed
to near-Gaussian (kurtosis $-0.09$) and reducing
L$^\infty$ norm by $40.1\%$. This is precisely
the concentration property the rotation is
designed to produce.

This result also validates the experimental design:
testing H1 on synthetic data alone would have
produced a misleading null result. Testing on real
KV vectors is the meaningful evaluation.

\subsection{Limitations}
\label{sec:limitations}

We identify three limitations that bound the
current claims:

\textbf{H2 Python implementation confound.}
Pipeline overhead is dominated by Python numpy
preprocessing at $d = 64$. Kyber-512 compiled C
($\approx 6~\mu$s) is faster than any Python-side
rotation implementation. This is an engineering
limitation, not an algorithmic one. A compiled
HD$^3$ kernel (C extension or CUDA) would reduce
preprocessing to $< 1~\mu$s, bringing total
pipeline latency within $\approx 2\times$ of
vanilla Kyber-512.

\textbf{Internal NTT not measurable.}
The \texttt{liboqs} \texttt{encap\_secret()} is a
black box. The hypothesis that reduced input
coefficient dynamic range decreases NTT conditional
reduction frequency cannot be verified with the
public API. An instrumented Kyber implementation
is required.

\textbf{Formal security reduction.}
The security argument is sound but informal.
We argue that pre-rotation does not alter the
Module-LWE problem instance and that Kyber's
IND-CCA2 proof holds for arbitrary message
distributions. A complete formal reduction
establishing composable security remains future work.

\subsection{Production Path}
\label{sec:production}

The practical deployment path for QuantRot-PQC
is straightforward:

\begin{enumerate}
    \item \textbf{Near-term:} Implement HD$^3$
    rotation as a C extension (e.g., via
    \texttt{ctypes} or \texttt{pybind11}).
    Expected preprocessing: $< 1~\mu$s at $d = 64$.
    Total pipeline: $\approx 7~\mu$s vs
    $6.25~\mu$s vanilla --- $\approx 12\%$ overhead.

    \item \textbf{Medium-term:} Integrate HD$^3$
    rotation inside an instrumented Kyber
    implementation to measure internal NTT
    reduction frequency directly.

    \item \textbf{Long-term:} Joint HD$^3$ + Kyber
    kernel that eliminates the preprocessing/
    encryption boundary entirely.
\end{enumerate}

\subsection{Combined Pipeline with TurboQuant}
\label{sec:combined}

A particularly efficient deployment combines
TurboQuant KV-cache compression with QuantRot-PQC
encryption. TurboQuant already applies HD$^3$
rotation as its first stage for near-Shannon-optimal
compression~\cite{turboquant2025}. In a system
already using TurboQuant, QuantRot-PQC adds
post-quantum security with \emph{zero additional
rotation cost} --- the rotation is shared between
the compression and encryption pipelines.

This combined pipeline simultaneously addresses
two orthogonal objectives: memory efficiency
(via TurboQuant compression) and transmission
security (via QuantRot-PQC encryption).

\subsection{Future Work}
\label{sec:future}

\begin{enumerate}
    \item Compiled HD$^3$ kernel to eliminate
    Python overhead and properly evaluate H2.
    \item Instrumented Kyber NTT to measure
    internal conditional reduction frequency
    before vs.\ after rotation.
    \item Formal IND-CCA2 security reduction
    for the QuantRot-PQC construction.
    \item Evaluation on larger models
    (Llama-3-8B, Gemma-7B) and higher
    dimensions ($d = 128$, $d = 256$).
    \item Generalization to other NIST PQC
    schemes (CRYSTALS-Dilithium, FALCON).
    \item Combined TurboQuant + QuantRot-PQC
    pipeline evaluation.
\end{enumerate}

%% ─────────────────────────────────────────────
\section{Conclusion}
\label{sec:conclusion}
We proposed QuantRot-PQC, a pre-encryption transformation
that applies HD$^3$ random unitary rotation to LLM
KV-cache vectors prior to CRYSTALS-Kyber polynomial
encoding. The key insight is that real transformer KV
vectors have heavy-tailed non-Gaussian distributions
--- unlike synthetic isotropic Gaussian vectors ---
and that HD$^3$ rotation concentrates this structure,
reducing L$^\infty$ norm and enabling more uniform
polynomial coefficient distributions before lattice
encoding.

Our main experimental result confirms Hypothesis H1:
HD$^3$ rotation reduces the L$^\infty$ norm of real
Llama-3.2-1B KV vectors by $40.1\%$ --- double the
$\geq 20\%$ threshold --- while kurtosis drops from
$6.15$ to $-0.09$, confirming near-Gaussian
concentration. The null result on synthetic isotropic
Gaussian vectors ($0.2\%$ reduction) is theoretically
consistent and serves as an experimental control:
rotation cannot concentrate what is already uniformly
distributed.

Pipeline latency (Hypothesis H2) is not supported
in the current Python implementation due to numpy
preprocessing overhead at $d = 64$. This is an
engineering limitation rather than an algorithmic
one: a compiled HD$^3$ kernel would reduce
preprocessing to $< 1~\mu$s, bringing total pipeline
latency within $\approx 2\times$ of vanilla
Kyber-512 ($6.25~\mu$s). Three open problems remain
for future work: a compiled HD$^3$ kernel, an
instrumented Kyber NTT to measure internal reduction
frequency, and a formal IND-CCA2 security reduction.

QuantRot-PQC addresses a timely and underexplored
problem: securing LLM inference traffic against
harvest-now-decrypt-later attacks while remaining
practical at production scale. The combination with
TurboQuant compression --- sharing the HD$^3$
rotation between compression and encryption at zero
additional rotation cost --- points toward a unified
efficient and secure LLM inference pipeline.

%% ─────────────────────────────────────────────
\bibliographystyle{alpha}
\begin{thebibliography}{99}

\bibitem{kyber2021}
R.~Avanzi et al.,
\textit{CRYSTALS-Kyber Algorithm Specifications and
Supporting Documentation}, Version 3.02, 2021.
\url{https://pq-crystals.org/kyber}

\bibitem{turboquant2025}
A.~Zandieh et al.,
\textit{TurboQuant: Online Vector Quantization with
Near-Optimal Distortion Rate},
arXiv:2504.19874, 2025.

\bibitem{brandao2016}
F.G.S.L.~Brand\~{a}o, A.W.~Harrow, M.~Horodecki,
\textit{Local Random Quantum Circuits are Approximate
Polynomial-Designs},
Communications in Mathematical Physics, 2016.
arXiv:1208.0692.

\bibitem{harrow2018}
A.W.~Harrow, S.~Mehraban,
\textit{Approximate Unitary $t$-Designs by Short Random
Quantum Circuits},
arXiv:1809.06957, 2018.

\bibitem{albrecht2015}
M.R.~Albrecht, R.~Player, S.~Scott,
\textit{On the Concrete Hardness of Learning with Errors},
Journal of Mathematical Cryptology, 9(3):169--203, 2015.

\bibitem{liboqs}
Open Quantum Safe Project,
\textit{liboqs: Open Quantum Safe Library},
\url{https://openquantumsafe.org}

\bibitem{llama2024}
Meta AI,
\textit{Llama 3.2: Lightweight Models},
\url{https://huggingface.co/meta-llama/Llama-3.2-1B},
2024.

\end{thebibliography}

\end{document}
